\textbf{Phần 1. Trắc nghiệm lựa chọn}
\setcounterpageref{ex}{0}

\begin{ex} %[3D14N70]
Kết quả của phép tính $20\,000 \times 4$ là
\choice
{$60000$}
{\True $80000$}
{$100000$}
{$70000$}
\loigiai{
$20000 \times 4 = 80000$.\\
\textbf{Đáp án của câu hỏi là B}
}
\end{ex}

\begin{ex} %[3D14N70]
Kết quả của phép nhân không nhớ $12\,321 \times 3$ là
\choice
{\True $36963$}
{$36953$}
{$36973$}
{$35963$}
\loigiai{
$12321 \times 3 = 36963$.\\
\textbf{Đáp án của câu hỏi là A}
}
\end{ex}

\begin{ex} %[3D14H70]
Kết quả của phép nhân có nhớ $15\,214 \times 4$ là
\choice
{$60846$}
{$60866$}
{$59856$}
{\True $60856$}
\loigiai{
$15214 \times 4 = 60856$.\\
\textbf{Đáp án của câu hỏi là D}
}
\end{ex}

\begin{ex} %[3D14H70]
Kết quả của phép nhân $24\,150 \times 3$ là
\choice
{$72440$}
{$72460$}
{\True $72450$}
{$71450$}
\loigiai{
$24150 \times 3 = 72450$.\\
\textbf{Đáp án của câu hỏi là C}
}
\end{ex}

\begin{ex} %[3D14H70]
Tìm $x$ biết: $x : 4 = 15\,000$.
\choice
{\True $60000$}
{$50000$}
{$70000$}
{$55000$}
\loigiai{
$x = 15000 \times 4 = 60000$.\\
\textbf{Đáp án của câu hỏi là A}
}
\end{ex}

\begin{ex} %[3D14H70]
Tính giá trị biểu thức: $14\,000 \times 5 + 20\,000$.
\choice
{$80000$}
{$100000$}
{\True $90000$}
{$85000$}
\loigiai{
$70000 + 20000 = 90000$.\\
\textbf{Đáp án của câu hỏi là C}
}
\end{ex}

\begin{ex} %[3D14V70]
Một nhà máy mỗi ngày sản xuất được $12\,500$ gói mì. Hỏi trong 6 ngày nhà máy đó sản xuất được bao nhiêu gói mì?
\choice
{$72500$ gói}
{\True $75000$ gói}
{$77500$ gói}
{$80000$ gói}
\loigiai{
$12500 \times 6 = 75000$ gói.\\
\textbf{Đáp án của câu hỏi là B}
}
\end{ex}

\begin{ex} %[3D14V70]
Một xe chở được $14\,200\text{ kg}$ hàng. Hỏi 5 xe như thế chở được bao nhiêu ki-lô-gam hàng?
\choice
{$70000\text{ kg}$}
{$72000\text{ kg}$}
{\True $71000\text{ kg}$}
{$75000\text{ kg}$}
\loigiai{
$14200 \times 5 = 71000\text{ kg}$.\\
\textbf{Đáp án của câu hỏi là C}
}
\end{ex}

\begin{ex} %[3D14V70]
Tính giá trị biểu thức: $18\,000 \times 4 - 32\,000$.
\choice
{\True $40000$}
{$38000$}
{$42000$}
{$45000$}
\loigiai{
$72000 - 32000 = 40000$.\\
\textbf{Đáp án của câu hỏi là A}
}
\end{ex}

\begin{ex} %[3D14V70]
Một hình vuông có cạnh dài $16\,000\text{ m}$. Chu vi hình vuông đó là
\choice
{$60000\text{ m}$}
{$68000\text{ m}$}
{$72000\text{ m}$}
{\True $64000\text{ m}$}
\loigiai{
$16000 \times 4 = 64000\text{ m}$.\\
\textbf{Đáp án của câu hỏi là D}
}
\end{ex}

\begin{ex} %[3D14V70]
Mua 3 bộ máy tính, mỗi bộ giá $18\,000\,000$ đồng (hoặc $18\,000$ nghìn đồng). Tổng số tiền mua máy tính là
\choice
{$52000$}
{\True $54000$}
{$56000$}
{$50000$}
\loigiai{
$18000 \times 3 = 54000$.\\
\textbf{Đáp án của câu hỏi là B}
}
\end{ex}

\begin{ex} %[3D14C70]
Một cửa hàng nhập về 4 bao đường, mỗi bao nặng $15\,200\text{ g}$. Cửa hàng đã bán được $35\,000\text{ g}$. Số đường còn lại là
\choice
{$25600\text{ g}$}
{$26000\text{ g}$}
{$25000\text{ g}$}
{\True $25800\text{ g}$}
\loigiai{
Tổng: $15200 \times 4 = 60800\text{ g}$.\\
Còn lại: $60800 - 35000 = 25800\text{ g}$.\\
\textbf{Đáp án của câu hỏi là D}
}
\end{ex}

\textbf{Phần 2. Trắc nghiệm đúng sai}
\setcounterpageref{ex}{0}

\begin{ex} %[3D14H70]
Xét tính đúng sai của các phép nhân:
\choiceTF
{\True $25000 \times 4 = 100000$.}
{\True $14200 \times 3 = 42600$.}
{$16000 \times 5 = 75000$.}
{\True $16000 \times 5 = 80000$.}
\loigiai{
\begin{itemize}
    \item $25000 \times 4 = 100000$: Đúng.
    \item $14200 \times 3 = 42600$: Đúng.
    \item $16000 \times 5 = 80000 \neq 75000$. Ý 3 Sai, Ý 4 Đúng.
\end{itemize}
\begin{itemize}
    \item a) \textbf{Đúng}.
    \item b) \textbf{Đúng}.
    \item c) \textbf{Sai}.
    \item d) \textbf{Đúng}.
\end{itemize}
}
\end{ex}

\begin{ex} %[3D14V70]
Một kho nhập về 3 xe hàng, mỗi xe chở $16\,500\text{ kg}$ gạo. Xét tính đúng sai:
\choiceTF
{\True Tổng số gạo 3 xe chở là $49\,500\text{ kg}$.}
{\True Nếu xuất đi $20\,000\text{ kg}$ thì còn lại $29\,500\text{ kg}$.}
{Nếu nhập thêm 1 xe như thế nữa thì tổng số gạo là $65\,000\text{ kg}$.}
{\True Nếu nhập thêm 1 xe như thế nữa thì tổng số gạo là $66\,000\text{ kg}$ ($16500 \times 4 = 66000$).}
\loigiai{
\begin{itemize}
    \item Tổng 3 xe: $16500 \times 3 = 49500\text{ kg}$. Đúng.
    \item Còn lại: $49500 - 20000 = 29500\text{ kg}$. Đúng.
    \item 4 xe chở: $16500 \times 4 = 66000\text{ kg} \neq 65000$. Ý 3 Sai, Ý 4 Đúng.
\end{itemize}
\begin{itemize}
    \item a) \textbf{Đúng}.
    \item b) \textbf{Đúng}.
    \item c) \textbf{Sai}.
    \item d) \textbf{Đúng}.
\end{itemize}
}
\end{ex}

\textbf{Phần 3. Trả lời ngắn}
\setcounterpageref{ex}{0}

\begin{ex} %[3D14H70]
Tính tích: $13\,210 \times 3$.
\shortans{39630}
\loigiai{
\textbf{Đáp số:} $39630$\\
$13210 \times 3 = 39630$.
}
\end{ex}

\begin{ex} %[3D14H70]
Tính tích: $24\,150 \times 4$.
\shortans{96600}
\loigiai{
\textbf{Đáp số:} $96600$\\
$24150 \times 4 = 96600$.
}
\end{ex}

\begin{ex} %[3D14V70]
Tính giá trị biểu thức: $15\,000 \times 5 - 25\,000$.
\shortans{50000}
\loigiai{
\textbf{Đáp số:} $50000$\\
$75000 - 25000 = 50000$.
}
\end{ex}

\begin{ex} %[3D14V70]
Một hình vuông có cạnh dài $18\,000\text{ m}$. Tính chu vi hình vuông theo mét.
\shortans{72000}
\loigiai{
\textbf{Đáp số:} $72000$\\
$18000 \times 4 = 72000\text{ m}$.
}
\end{ex}

\begin{ex} %[3D14V70]
Tìm $x$ biết: $x : 6 = 14\,000$.
\shortans{84000}
\loigiai{
\textbf{Đáp số:} $84000$\\
$x = 14000 \times 6 = 84000$.
}
\end{ex}

\begin{ex} %[3D14C70]
Một hình chữ nhật có chiều rộng $12\,500\text{ m}$, chiều dài gấp 3 lần chiều rộng. Tính chiều dài hình chữ nhật theo mét.
\shortans{37500}
\loigiai{
\textbf{Đáp số:} $37500$\\
$12500 \times 3 = 37500\text{ m}$.
}
\end{ex}
